\documentclass[12pt]{article}
\usepackage[utf8]{inputenc}
\usepackage[spanish]{babel}
\usepackage{amsmath}
\usepackage{graphicx}
\usepackage[margin=1in]{geometry} % Para controlar los márgenes
\usepackage{amsfonts}
\usepackage{tikz}
\begin{document}

% --- CARÁTULA ---
\begin{titlepage}
    \centering
    {\bfseries\LARGE UNIVERSIDAD NACIONAL MAYOR DE SAN MARCOS \par}
    {\scshape\Large (UNIVERSIDAD DEL PERÚ, DECANA DE AMÉRICA) \par}
    \vspace{1cm}
    
    \includegraphics[width=0.3\textwidth]{logo_unmsm}\par\vspace{0.5cm}
    
    {\bfseries\large FACULTAD DE CIENCIAS ECONÓMICAS \par}
    \vspace{1cm}
    
    {\bfseries\Large TAREA N°1: \par}
    {\Large "Price competition with uncertain costs" \par}
    \vspace{1cm}
    
    {\bfseries\large Grupo 4 \par}
    {\bfseries\large Integrantes \par}
    \vspace{0.5cm}
    CARO SANDONAS, JOSE GUILLERMO \\
    DURAND APAZA, YADIRA NICOL \\
    CHUMACERO PEÑA LEANDRO JOSUET \\
    QUISPE SANTOYO, SEBASTIAN GABRIEL \\
    APELLIDO9 APELLIDO10 \\
    \vspace{0.5cm}
    
    {\bfseries\large Docente: \par}
    Huari Leasaski, David \\
    \vfill
     \vspace{0.5cm}
    {\large Lima, Perú \par}
        \vspace{0.5cm}
    {\large 29 de abril de 2025 \par}
\end{titlepage}

% --- ÍNDICE ---
\newpage
\tableofcontents
\newpage

% --- CONTENIDO ---
\section{Introducción}
Este trabajo analiza el modelo de competencia de precios bajo condiciones de incertidumbre en los costos...
\subsection{Contraste de Modelos: De la Información Perfecta a la Asimetría,(\textit{Trade-off}) y Equilibrio de Nash Bayesiano}
En el modelo de Bertrand estándar (información perfecta), todas las empresas conocen exactamente la estructura de costos de sus competidores. Esta dinámica genera una competencia feroz que empuja los precios hasta igualar el costo marginal ($p = c$), haciendo que las firmas solo cubran costos sin generar beneficios. Por el contrario, en el modelo con información asimétrica (costos inciertos), el costo $c_i$ de cada empresa es información privada o secreta''. Las firmas rivales no conocen el valor exacto, solo la distribución estadística de dichos costos. Esta opacidad'' permite a las empresas cobrar un margen de beneficio (\textit{markup}) positivo, generando lo que se conoce como ``rentas de eficiencia''.
Esta asimetría transforma la naturaleza de la competencia. Mientras que en el modelo clásico reducir el precio en un centavo garantiza capturar todo el mercado, bajo incertidumbre la empresa enfrenta una disyuntiva constante: fijar un precio más alto para extraer mayores rentas de información aumenta su margen de ganancia ($p_i - c_i$), pero disminuye drásticamente la probabilidad de que su oferta sea la más baja, elevando el riesgo de perder la totalidad de la demanda frente a un rival estadísticamente más eficiente.
Desde una perspectiva metodológica, la resolución de este problema requiere ir más allá de un equilibrio tradicional. El resultado de este modelo es un Equilibrio de Nash Bayesiano. Al desconocer los tipos'' exactos (costos) de sus rivales, cada firma debe formular creencias'' probabilísticas (en este caso, asumiendo una distribución uniforme $U \sim [0,1]$) para lograr maximizar su beneficio esperado balanceando la disyuntiva entre el margen y el volumen de ventas.

Como demuestra el artículo original de Spulber (1995), esta dinámica resuelve la Paradoja de Bertrand. El mercado ya no salta abruptamente de un escenario de poder de mercado a uno de ganancias nulas. En su lugar, el modelo evidencia una transición continua hacia la eficiencia competitiva: a medida que el número de firmas ($n$) se incrementa, la asimetría se diluye y el margen óptimo de cada jugador se reduce progresivamente.

\section{Desarrollo del Modelo}
\subsection{Planteamiento y demanda}

El primer paso para analizar la competencia de precios bajo incertidumbre consiste en definir la estructura de la demanda que enfrenta cada firma. Dado que las empresas venden un bien homogéneo, los consumidores basarán su decisión de compra exclusivamente en el precio más bajo disponible en el mercado.

Desde la perspectiva de la empresa $i$, se deben distinguir tres escenarios posibles para su demanda individual ($q_i$), dependiendo de la relación entre su precio ($p_i$) y el precio mínimo de sus competidores ($\hat{p}_{-i}$):

\begin{enumerate}
    \item \textbf{Liderazgo en precios ($p_i < \hat{p}_{-i}$):} Cuando la empresa $i$ fija un precio estrictamente inferior al de todos sus competidores, captura la totalidad de la demanda del mercado, definida por la función lineal $q_i = 1 - p_i$.
    
    \item \textbf{Exclusión del mercado ($p_i > \hat{p}_{-i}$):} Si el precio de la empresa $i$ supera al menos al de uno de sus competidores, la cantidad demandada es nula ($q_i = 0$), ya que los consumidores optarán por la opción más económica.
    
    \item \textbf{Empate en el precio mínimo ($p_i = \hat{p}_{-i}$):} En caso de que el precio de la empresa coincida con el mínimo de la competencia, la demanda se distribuye de manera equitativa entre las $m$ firmas que comparten dicho precio, resultando en $q_i = \frac{1 - p_i}{m}$.
\end{enumerate}

\subsubsection{Representación Formal}

Matemáticamente, la función de demanda de la empresa $i$ se expresa mediante la siguiente función por partes:

\begin{equation}
    q_{i}(p_{i}, \hat{p}_{-i}) = 
    \begin{cases} 
        1 - p_{i} & \text{si } p_{i} < \hat{p}_{-i} \\ 
        \frac{1 - p_{i}}{m} & \text{si } p_{i} = \hat{p}_{-i} \\ 
        0 & \text{en cualquier otro caso} 
    \end{cases}
\end{equation}


\subsection{Maximización de Beneficios y Equilibrio de Nash Bayesiano}

En esta sección se analiza cómo las firmas toman decisiones cuando no conocen con certeza los costos de sus competidores. Para resolver este dilema, utilizaremos el concepto de \textbf{Equilibrio de Nash Bayesiano (BNE)}, que es la herramienta estándar para estudiar juegos con información privada \cite{harsanyi}.

\subsubsection{Supuestos Fundamentales}
\begin{itemize}
    \item \textbf{Estrategia Óptima ($p^*$):} Representamos como $p^*(c)$ a la función de precios que una firma elige racionalmente según su costo marginal.
    \item \textbf{Estrategia Monótona:} Se asume que $p^*(c)$ es estrictamente creciente. Esto es fundamental para que exista su función inversa ($p^{*-1}$), la cual permite realizar el proceso inverso: conocer el costo a partir de un precio dado.
    \item \textbf{Distribución de Costos:} Se asume que los costos marginales $c_i$ se distribuyen de forma uniforme entre 0 y 1 ($U[0,1]$).
\end{itemize}

\subsubsection{Derivación de la Probabilidad de Éxito}
Para que la empresa $i$ capture el mercado, el precio que ella elige ($p_i$) debe ser menor al precio de equilibrio que elijan sus competidores ($p^*(c_j)$). La probabilidad frente a $n-1$ competidores independientes es:
\begin{equation}
    Pr(p_i < \hat{p}_{-i}) = Pr(p_i < p^*(c_1)) \times Pr(p_i < p^*(c_2)) \times \dots \times Pr(p_i < p^*(c_{n-1}))
\end{equation}

\noindent \textbf{Transformación paso a paso} \\ 
Para calcular esta probabilidad, aplicamos la función inversa $p^{*-1}$ a ambos lados de la desigualdad interna. Este proceso permite despejar el costo del competidor ($c_j$):

\begin{enumerate}
    \item \textbf{Estado inicial:} Comparación de precios.
    \begin{equation} p_i < p^*(c_j) \end{equation}
    \item \textbf{Aplicación de la inversa en ambos lados:}
    \begin{equation} p^{*-1}(p_i) < p^{*-1}(p^*(c_j)) \end{equation}
    \item \textbf{Cancelación matemática:} Dado que una función aplicada a su propia inversa se anula la expresión se simplifica a una comparación de costos:
    \begin{equation} p^{*-1}(p_i) < c_j \end{equation}
\end{enumerate}

\noindent \textbf{Resultado probabilístico} \\
Bajo la distribución uniforme $U[0,1]$, la probabilidad de que el costo real del rival ($c_j$) sea mayor que el valor $p^{*-1}(p_i)$ es simplemente la diferencia hasta el límite superior:
\begin{equation}
    Pr(c_j > p^{*-1}(p_i)) = 1 - p^{*-1}(p_i)
\end{equation}

\subsubsection{Planteamiento Final del Beneficio Esperado}
Al combinar la probabilidad acumulada de los $n-1$ rivales con el margen y la demanda, el problema de maximización queda definido como:

\begin{equation}
    \max_{p_{i} \ge 0} E[\pi_i] = \underbrace{(p_{i} - c_{i})}_{\text{Margen}} \times \underbrace{(1 - p_{i})}_{\text{Demanda}} \times \underbrace{[1 - p^{*-1}(p_{i})]^{n-1}}_{\text{Prob. de ganar}}
\end{equation}
\subsection{Derivamos el problema de maximización del beneficio esperado con respecto al precio p_i}

\subsubsection{Condición de primer orden} 
En este apartado se deriva la condición de primer orden del problema de maximización de beneficios. El objetivo es mostrar paso a paso cómo se obtiene la expresión que caracteriza el precio óptimo, detallando el uso de las reglas de derivación y su interpretación económica.
Definimos los tres bloques que conforman la función de beneficio esperado:
\begin{itemize}
    \item \textbf{Bloque 1 (Margen):} $U = (p_i - c_i)$ \quad $\implies$ \quad $U' = 1$
    \item \textbf{Bloque 2 (Demanda):} $V = (1 - p_i)$ \quad $\implies$ \quad $V' = -1$
    \item \textbf{Bloque 3 (Probabilidad):} $W = [1 - p^{*-1}(p_i)]^{n-1}$ 
\end{itemize}

Para el Bloque 3, aplicamos la regla de la cadena, obteniendo su derivada ($W'$):
\begin{equation}
    W' = (n-1) [1 - p^{*-1}(p_i)]^{n-2} \cdot \left( - \frac{d p^{*-1}(p_i)}{d p_i} \right)
\end{equation}

\noindent \textbf{Aplicación de la Regla del Producto} \\
La derivada total de la función $E[\pi_i] = U \cdot V \cdot W$ se construye sumando el efecto de cada bloque:
\begin{equation}
    \frac{d E[\pi_i]}{d p_i} = \underbrace{U'VW}_{\text{Efecto Margen}} + \underbrace{UV'W}_{\text{Efecto Demanda}} + \underbrace{UVW'}_{\text{Efecto Probabilidad}} = 0
\end{equation}

Sustituyendo los valores, obtenemos la ecuación diferencial base:
\begin{equation}
    (1-p_i)[1-p^{*-1}]^{n-1} - (p_i-c_i)[1-p^{*-1}]^{n-1} - (p_i-c_i)(1-p_i)(n-1)[1-p^{*-1}]^{n-2} \frac{d p^{*-1}}{d p_i} = 0
\end{equation}

\subsubsection{Simplificación y Estructura Final}

Para resolver la ecuación diferencial (10), observamos que el término de probabilidad $[1 - p^{*-1}(p_i)]$ se repite en toda la expresión. Al simplificar los términos comunes y reordenar los componentes para aislar la derivada de la función inversa, llegamos a la representación técnica de la CPO:

\begin{equation}
    \frac{(1 - 2p_i + c_i)(1 - p^{*-1}(p_i))}{(n - 1)(p_i - c_i)(1 - p_i)} = \frac{d p^{*-1}(p_i)}{d p_i}
\end{equation}

\noindent Esta condición final resume el dilema fundamental de la firma, el cual ha sido ampliamente documentado en la literatura de Organización Industrial \cite{tirole1988}. Como se observa en la ecuación (11), el equilibrio se alcanza al balancear dos fuerzas económicas críticas:

\begin{itemize}
    \item \textbf{Efecto Margen:} Un aumento en el precio eleva la utilidad por unidad vendida.
    \item \textbf{Efecto Probabilidad:} Siguiendo la lógica de las subastas de primer precio, un precio más alto reduce la probabilidad de que la firma sea la opción más barata del mercado \cite{klemperer1999}.
\end{itemize}

\subsection{Función de precios lineal} 
En esta sección se resuelve la ecuación diferencial obtenida en la CPO para encontrar la función de precios de equilibrio. Para ello, se propone una forma funcional y se determinan sus coeficientes mediante el método de igualación.

\subsubsection{Propuesta de Forma Funcional}
Asumimos que la estrategia de equilibrio para cada firma es una función lineal de su costo marginal:
\begin{equation}
    p^*(c_i) = a + b c_i
\end{equation}

Donde $a$ representa el componente fijo o margen base (\textit{markup}) y $b$ representa la sensibilidad del precio ante cambios en el costo.

\subsubsection{Derivación de la Función Inversa}
Dado que la condición de primer orden (11) requiere la función inversa $p^{*-1}(p_i)$ y su derivada, procedemos a despejar el costo marginal ($c_i$) de la propuesta lineal anterior:
\begin{equation}
    c_i = p^{*-1}(p_i) = \frac{p_i - a}{b}
\end{equation}

A partir de este despeje, calculamos la derivada de la función inversa con respecto al precio, la cual es constante en el caso lineal:
\begin{equation}
    \frac{d p^{*-1}(p_i)}{d p_i} = \frac{1}{b}
\end{equation}

\subsubsection{Resolución de coeficientes y solución final}

Para simplificar la CPO, primero expresamos los términos de la ecuación (11) en función de nuestra propuesta lineal $p = a + bc_i$. Partiendo de la condición de que los pesos de la estrategia deben sumar la unidad ($a + b = 1 \implies b = 1 - a$), derivamos las siguientes identidades:

\begin{itemize}
    \item \textbf{Término de margen ($1-p$):}
    \begin{equation}
        1 - p = 1 - (a + bc_i) = (1 - a) - (1 - a)c_i = (1 - a)(1 - c_i)
    \end{equation}
    
    \item \textbf{Término de diferencial de precios ($p-c_i$):}
    \begin{equation}
        p - c_i = (a + bc_i) - c_i = a - (1 - b)c_i = a - ac_i = a(1 - c_i)
    \end{equation}
    
    \item \textbf{Término del numerador ($1 + c_i - 2p$):}
    \begin{equation}
        1 + c_i - 2(a + bc_i) = (1 - 2a) + c_i(1 - 2b) = (1 - 2a)(1 - c_i)
    \end{equation}
\end{itemize}

\subsubsection{Sustitución y Resolución Algebraica}

Ahora, sustituimos estas identidades y la derivada (14) en la CPO (11). El reemplazo se visualiza de la siguiente forma:

\begin{equation}
    \frac{\overbrace{(1 - 2a)(1 - c_i)}^{1 - 2p_i + c_i} \cdot \overbrace{(1 - c_i)}^{1 - p^{*-1}}}{(n - 1) \cdot \underbrace{a(1 - c_i)}_{p_i - c_i} \cdot \underbrace{(1 - a)(1 - c_i)}_{1 - p_i}} = \frac{1}{b}
\end{equation}

Al observar la fracción, notamos que el término $(1 - c_i)$ aparece multiplicando dos veces tanto en el numerador como en el denominador, permitiendo su cancelación total. Simplificando y recordando que $(1 - a) = b$, la ecuación se reduce a:

\begin{equation}
    \frac{1 - 2a}{(n - 1)a \cdot b} = \frac{1}{b}
\end{equation}

Multiplicando ambos miembros por $b$, obtenemos la ecuación lineal:
\begin{equation}
    1 - 2a = (n - 1)a \implies 1 = na + a \implies a = \frac{1}{n+1}
\end{equation}

Finalmente, dado que $b = 1 - a$, obtenemos el valor del segundo coeficiente:
\begin{equation}
    b = 1 - \frac{1}{n+1} = \frac{n}{n+1}
\end{equation}
\section{Interpretación y  análisis}
\subsection{Interpretación y Discusión: El Rol de la Información y la Competencia}

A partir de la función de equilibrio obtenida, $p^*(c_i) = \frac{1}{n+1} + \frac{n}{n+1}c_i$, es posible extraer conclusiones fundamentales sobre el comportamiento estratégico de las firmas bajo condiciones de incertidumbre.

\subsection{Rentas de Información y Markup}
\begin{itemize}
    \item El modelo demuestra que cada firma fija un precio estrictamente superior a su costo marginal para todo $0 \le c_i < 1$.
    \item Esta diferencia constituye una \textbf{renta de información}: la firma capitaliza el hecho de que sus rivales desconocen su nivel real de eficiencia para extraer un excedente adicional.
    \item Existe un balance estratégico entre aumentar el margen de ganancia por unidad y reducir la probabilidad de ser el postor más bajo del mercado.
\end{itemize}

\subsection{Efecto del Número de Firmas ($n$)}
La estructura de la solución permite observar la sensibilidad del equilibrio ante cambios en la estructura del mercado:
\begin{itemize}
    \item \textbf{Presión Competitiva:} A medida que $n$ aumenta, el término constante $\frac{1}{n+1}$ disminuye, lo que reduce el \textit{markup} base de las firmas.
    \item \textbf{Convergencia de Bertrand:} En el límite, cuando $n \to \infty$, el precio converge al costo marginal ($p \to c_i$).
    \item Esta convergencia implica que la incertidumbre pierde relevancia estratégica frente a la presión competitiva de un gran número de rivales, recuperando el resultado del modelo de Bertrand clásico.
\end{itemize}

\subsection{Justificación de la Condición de Borde ($p(1)=1$)}
La fijación del precio en el extremo superior de la distribución de costos responde a una lógica de equilibrio de Nash[cite: 1]:
\begin{itemize}
    \item \textbf{Racionalidad Económica:} Una firma con $c_i = 1$ no puede cobrar un precio menor a 1 sin incurrir en pérdidas operativas directas[cite: 1].
    \item \textbf{Inviabilidad de Markup:} Un precio superior a 1 garantiza ventas nulas, ya que la probabilidad de que todos los demás rivales tengan también el costo máximo es negligible en una distribución continua[cite: 1].
\end{itemize}
\section{Aplicación Empírica: Licitaciones Públicas y Subastas a Sobre Cerrado}

El marco teórico de la competencia en precios con costos inciertos encuentra su aplicación empírica más directa en los procesos de licitaciones públicas (subastas a sobre cerrado) y contratos de suministro corporativo (B2B). En estos escenarios, el ``bien homogéneo'' está rígidamente definido por los términos de referencia o expedientes técnicos (por ejemplo, la construcción de una infraestructura pública con especificaciones estandarizadas). Las firmas participantes conocen con precisión sus propios costos operativos ($c_i$), pero la estructura de costos de sus rivales constituye información estrictamente privada. 

Al formular su propuesta económica, cada postor resuelve en la práctica un Equilibrio de Nash Bayesiano: debe decidir qué margen de ganancia agregar a su costo base, sabiendo que un margen elevado incrementa la rentabilidad potencial del proyecto, pero penaliza drásticamente la probabilidad de adjudicarse el contrato. Asimismo, la dinámica poblacional del modelo se evidencia empíricamente: en licitaciones con escasa participación de postores (un $n$ reducido), el Estado o comprador corporativo suele pagar primas más altas (\textit{markups}); por el contrario, en procesos con alta concurrencia (un $n$ elevado), el miedo a perder la adjudicación obliga a las firmas a comprimir sus márgenes, haciendo que las ofertas converjan hacia su costo marginal y maximizando la eficiencia en el uso de los recursos públicos o privados.

\section{Análisis Comparativo y Conclusiones Económicas}

\subsection{Tabla Comparativa: Efecto de la Competencia sobre el Precio}

Para ilustrar la convergencia del precio hacia el costo marginal, evaluamos la función de equilibrio $p^*(c_i) = \frac{1}{n+1} + \frac{n}{n+1}c_i$ bajo distintos escenarios de competencia ($n$) y niveles de eficiencia ($c_i$):

\begin{table}[h]
\centering
\begin{tabular}{|c|c|c|c|c|}
\hline
\textbf{Costo ($c_i$)} & \textbf{n = 2} (Duopolio) & \textbf{n = 5} & \textbf{n = 10} & \textbf{n $\to \infty$} (Bertrand) \\ \hline
0.0 & 0.33 & 0.17 & 0.09 & 0.00 \\ \hline
0.2 & 0.47 & 0.33 & 0.27 & 0.20 \\ \hline
0.5 & 0.67 & 0.58 & 0.55 & 0.50 \\ \hline
0.8 & 0.87 & 0.83 & 0.82 & 0.80 \\ \hline
1.0 & 1.00 & 1.00 & 1.00 & 1.00 \\ \hline

\end{tabular}


\caption{Evolución del precio de equilibrio según el número de firmas y el costo marginal.}
\end{table}

Como se observa en la tabla, a medida que el número de firmas ($n$) aumenta, el precio de equilibrio para cualquier nivel de costo se reduce progresivamente, acercándose a la línea de 45 grados donde $p = c_i$.

\subsection{Analisis gráfico: Efecto de la competencia sobre el precio y los márgenes de ganacia (markup)}
\begin{itemize}
    \item Gráficamente, la función de precios $p(c_{i})$ será siempre una línea recta continua, con variaciones en su ángulo de inclinación dependientes de manera exclusiva del valor de sus parámetros estructurales, específicamente del número de firmas $n$. La función tiene su origen en el intercepto:
   
    $$ p(0) = \frac{1}{n+1} $$
   
    Esto ocurre en el escenario límite donde el costo marginal de la firma $i$ alcanza su cota inferior ($c_{i}=0$). A partir de este punto, su valor se incrementa monotónicamente conforme la firma se vuelve menos eficiente (es decir, conforme crecen los costos $c_i$) a una tasa de cambio constante equivalente a $\frac{n}{n+1}$.

    \item En términos geométricos y analíticos, la función de precios $p(c_{i})$ se encuentra generalmente por encima de la línea de $45^\circ$ del plano cartesiano (la cual representa la función identidad $p = c_i$), especialmente cuando la estructura de costos de la firma refleja una alta eficiencia relativa (costos marginales $c_i$ tendientes a cero). Se distinguen de manera rigurosa tres situaciones al respecto de esta relación:
   
    \begin{itemize}
        \item \textbf{Situación I (Eficiencia y margen positivo):} Las firmas fijan un precio estrictamente superior a su costo marginal, garantizando un margen de ganancia unitario positivo. Esto se cumple para todo el intervalo definido por la inecuación:
       
        $$ p(c_{i}) > c_{i} \quad \forall \ c_{i} \in [0, 1) $$
       
        \item \textbf{Situación II (Límite competitivo de rentabilidad nula):} Se establece un precio que coincide exactamente con su costo marginal, lo que implica un mark-up o margen de ganancia igual a cero. Este escenario ocurre puntualmente en el extremo superior del soporte de la distribución de costos:
       
        $$ p(c_{i}) = c_{i} \quad \text{cuando} \ c_{i} = 1 $$
       
        \item \textbf{Situación III (No viabilidad teórica de operaciones):} Ocurre una situación de inviabilidad económica de las operaciones si el costo llegara a superar la unidad. En esta región paramétrica, la función de precio óptima se ubicaría por debajo de la curva de costos marginales:
       
        $$ p(c_{i}) < c_{i} \quad \forall \ c_{i} \in (1, \infty) $$
       
        No obstante, dado que se ha asumido una distribución de costos uniforme $c_i \sim U[0,1]$, este tercer escenario queda estrictamente fuera del soporte del modelo probabilístico planteado.
    \end{itemize}
\end{itemize}

\vspace{1cm}

\begin{figure}[htbp] % h: aquí, t: arriba, b: abajo, p: página especial
    \centering
    
    \begin{tikzpicture}[scale=3.5]
        % Plano cartesiano completo
        % Eje X (costos)
        \draw[<->, thick] (-0.3, 0) -- (1.6, 0) node[right] {$c_i$ (Costo Marginal)};
        % Eje Y (precios)
        \draw[<->, thick] (0, -0.3) -- (0, 1.6) node[above] {$p(c_i)$ (Precio)};
    
        % Cuadrícula de referencia
        \draw[step=0.5, gray!30, very thin] (-0.2, -0.2) grid (1.5, 1.5);
    
        % Etiqueta del origen
        \node[below left] at (0,0) {$0$};
    
        % Línea de 45 grados (Función Identidad)
        \draw[dashed, blue, thick, domain=-0.2:1.4] plot (\x, \x) node[above right] {$p = c_i$};
    
        % Función de precio p(c_i) asumiendo n=2
        \draw[red, thick, domain=-0.2:1.4] plot (\x, {0.333 + 0.666*\x}) node[right] {$p(c_i) = \frac{1}{n+1} + \frac{n}{n+1}c_i$};
    
        % Puntos de intersección y marcadores
        \node[left, fill=white, inner sep=1pt] at (0, 0.333) {$\frac{1}{n+1}$};
        \draw[fill=black] (0, 0.333) circle [radius=0.015];
    
        % Proyecciones ortogonales para c_i = 1
        \draw[dotted, thick] (1, 0) node[below, fill=white, inner sep=1pt] {$1$} -- (1, 1) -- (0, 1) node[left, fill=white, inner sep=1pt] {$1$};
        \draw[fill=black] (1, 1) circle [radius=0.015];
        
        % Marcas en los ejes
        \draw[thick] (-0.1, 0.02) -- (-0.1, -0.02);
        \draw[thick] (0.02, -0.1) -- (-0.02, -0.1);
    
        % Etiquetas de las tres situaciones
        \node[align=center, font=\footnotesize, text width=2.8cm, blue!80!black, fill=white, fill opacity=0.85, text opacity=1] at (0.45, 0.95) {\textbf{Situación I}\\ $p(c_i) > c_i$ \\ $c_i \in [0, 1)$};
        
        % Flecha y etiqueta para Situación II
        \draw[->, shorten >=2pt, thick] (1.3, 0.6) -- (1.02, 0.98);
        \node[align=center, font=\footnotesize, text width=2cm, fill=white, fill opacity=0.85, text opacity=1] at (1.3, 0.45) {\textbf{Situación II}\\ $p(c_i) = c_i$ \\ $c_i = 1$};
        
        % Etiqueta para la Situación III
        \node[align=center, font=\footnotesize, text width=2.5cm, gray, fill=white, fill opacity=0.85, text opacity=1] at (1.3, 1.35) {\textbf{Situación III}\\ $p(c_i) < c_i$ \\ $c_i > 1$};
    \end{tikzpicture}

    \vspace{0.3cm} % Espacio ajustable entre el dibujo y el título
    \caption{Relación analítica entre el costo marginal y la función de precios óptima.}
    \label{fig:analisis_precios}
\end{figure}

\subsection{Conclusión Económica: La Incertidumbre como Escudo}

A diferencia del modelo de \textbf{Bertrand clásico} con información perfecta, donde la presencia de solo dos empresas es suficiente para que el precio colapse instantáneamente al costo marginal ($p = MC$), este modelo revela una dinámica de caída más lenta y gradual. 

La \textbf{incertidumbre sobre los costos} de los rivales actúa como un "escudo" o colchón estratégico para las firmas. Al no saber con certeza si el competidor es más eficiente, las empresas mantienen la capacidad de fijar precios por encima de sus costos marginales sin temor a perder la totalidad del mercado de manera inmediata[cite: 1]. Este fenómeno explica por qué, en mercados reales con información asimétrica (como licitaciones o sectores industriales complejos), los márgenes de ganancia no desaparecen tan rápido como predice la teoría tradicional, incluso ante un incremento en el número de competidores.
%Bibliografía
\begin{thebibliography}{5}

\bibitem{harsanyi1967} 
Harsanyi, J. C. (1967). Games with incomplete information played by ``Bayesian'' players. \textit{Management Science}, 14(3), 159-182.

\bibitem{tirole1988} 
Tirole, J. (1988). \textit{The Theory of Industrial Organization}. MIT Press.

\bibitem{munoz2026} 
Munoz-Garcia, F. (2026). \textit{Price Competition with Uncertain Costs}. EconS 594 - Industrial Organization. Washington State University.

\bibitem{spulber1995} 
Spulber, D. F. (1995). Bertrand Competition when Costs are Unknown. \textit{The Journal of Industrial Economics}, 43(1), 1-11.

\bibitem{klemperer1999} 
Klemperer, P. (1999). Auction Theory: A Guide to the Literature. \textit{Journal of Economic Surveys}, 13(3), 227-286.

\end{thebibliography}
\end{document}